\section{质量管理}

\subsection{质量目标}

本项目以“高可靠性、高可用性、高可维护性”为核心质量目标，具体包括：
\begin{itemize}
	\item \textbf{功能正确性：} 所有功能模块均严格按照需求文档实现，确保业务流程无误，数据处理精准。
	\item \textbf{系统稳定性：} 通过全面测试与异常处理，保障系统在高并发、复杂场景下依然稳定运行。
	\item \textbf{可维护性：} 采用模块化设计与规范化编码，便于后续扩展与维护。
	\item \textbf{可用性：} 前端界面友好，交互流畅，用户体验优异，满足实际教学与科研需求。
\end{itemize}

\subsection{评审机制}

为确保项目质量，团队建立了全流程、多层次的评审机制：
\begin{itemize}
	\item \textbf{需求评审：} 产品经理组织全员参与，确保需求完整、无歧义，提前规避实现风险。
	\item \textbf{设计评审：} 由后端、前端及数据负责人共同参与，评估系统架构、接口设计与数据流合理性。
	\item \textbf{代码评审：} 依托 GitHub Pull Request 机制，所有核心代码均需经至少两人 Code Review，提升代码质量与可读性。
	\item \textbf{测试评审：} 测试负责人定期组织测试用例与结果评审，确保测试覆盖全面、缺陷及时闭环。
\end{itemize}

\subsection{测试计划}

测试工作贯穿项目全周期，覆盖单元、集成与验收各环节，责任到人，保障系统交付高质量。

\begin{longtable}{|p{32mm}|p{72mm}|p{44mm}|}
\hline
\textbf{测试类型} & \textbf{测试内容} & \textbf{负责人} \\
\hline
单元测试 & 核心函数、组件和服务逻辑 & 杨吉祥、田红瑾 \\
\hline
集成测试 & 模块间接口、数据流和异常处理 & 杨吉祥、田红瑾 \\
\hline
验收测试 & 主要业务流程和课程验收场景 & 杨吉祥、田红瑾 \\
\hline
\end{longtable}

通过上述严密的质量管理体系，团队能够有效防控风险、持续优化产品，确保 Minerva 系统以稳定质量顺利交付。

